\section{Discussion}

\textit{Measure overlap before claiming generalization.} The clearest finding
here is about the benchmarks rather than the models. A paper that trains on the
Kaggle aggregate and tests on SpamAssassin is testing mostly on its own
training data, and exact-match deduplication does not reveal it because the
aggregate changed the line breaks. A near-duplicate pass over 170{,}000 emails
took five minutes on a laptop in our run, so there is no practical reason to
skip it before a cross-dataset claim. Our random-removal control also shows
that the effect is leakage and not the smaller training set, which is the first
question a careful reader asks.

\textit{In-corpus scores answer a different question.} The 97 to 99\% in the
literature is real, but it says how well a model recognises more mail from the
same collection. A mail server does not receive mail from the training
collection. On unseen corpora the same model loses several points; on
AI-written mail it loses many more; and the model with the best in-corpus score
in our table is among the weakest on AI-written phishing.

\textit{Know what the positive class is.} Only a fifth of the positive emails
in these corpora are phishing. Work that reports "phishing detection" on them
is largely reporting spam detection. This does not invalidate the corpora, but
a paper should either say which it means or, as we did, measure the split.

\textit{Choose an operating point, not a leaderboard position.} At the balanced
prevalence used in benchmarks, several models look similar. At a realistic
phishing rate the false alarm rate dominates, and models that share an F1
differ by an order of magnitude in how much legitimate mail they would flag.
The two-stage detectors show a more practical way to use a cheap model: a free
classical filter in front of an LLM cuts both the cost and the false alarms,
and the reverse order is the right shape for a quarantine queue that a person
reviews. Conformal deferral~\cite{triage} is the natural next step.

\textit{What we would deploy.} For an organisation with no GPU and a small
budget, a TF-IDF filter in front of zero-shot Qwen-2.5-7B is the best
combination we measured: it is the cheapest configuration in our tables and has
the lowest false alarm rate. If the budget allows a commercial API,
Gemini-3.1-Flash-Lite was better than every open model on every axis we
measured, which is an honest result even though it is not the one we hoped for.
Gemma-3-12B should be used only where a person reviews the output.

\textit{Test the shortcut you suspect, on the same emails.} The AI-written
corpus writes links as a placeholder that appears almost only in phishing mail,
and the obvious comparison, splitting the phishing emails by whether they carry
it, says the token is worth a large amount of recall. That comparison is
confounded, because the two halves are different kinds of attack. Removing the
token from the same emails moves the verdicts by almost nothing. We report both
because the quick test pointed the wrong way, and a reader who only ran it
would have drawn the opposite conclusion about this corpus.

\section{Limitations}

\begin{itemize}
\item \textit{Labels.} Four corpora count ordinary spam as positive. We measure
the split with an LLM annotator and a second annotator for agreement, but the
underlying labels stay as published, and the annotator is itself a model.
\item \textit{Sample size.} Each test set has 300 emails, 150 for the one-class
sets. The bootstrap interval on the mean F1 is about $\pm$0.015; the
leave-one-corpus-out jackknife is wider and is the number to read when asking
about a new corpus.
\item \textit{Input budgets differ by model family.} TF-IDF sees 20{,}000
characters, the LLMs 1{,}500, DistilBERT 128 word pieces. DistilBERT is
therefore the most truncated model and its row should be read as a lower bound.
\item \textit{Training budgets differ.} DistilBERT trains on about 5{,}000
emails per fold against 40{,}000 for the classical models, to fit the laptop.
\item \textit{LLM pre-training.} The legacy corpora may be inside the LLMs'
pre-training data~\cite{contamsurvey}. E-PhishLLM was released in 2025, after
most of these models, which makes it the most trustworthy of our tests, and the
published detectors we ran are contaminated by construction.
\item \textit{Providers.} OpenRouter routes each model to a provider of its
choosing, so a small part of the cost and a small part of the behaviour is
routing rather than the model.
\item \textit{Scope.} Text only, one prompt per model, one run. Headers, URLs
and attachments are ignored, and the Kaggle corpus has no subject line at all,
which is an unmeasured difference between it and the others.
\item \textit{Phishsense-1B} is evaluated through an ungated third-party merge
and with our own prompt, because the original repository and its prompt
template are gated. Its numbers may understate the released model.
\end{itemize}

\section{Conclusion}

Public phishing corpora overlap heavily, and in a way that exact matching
cannot see: most of the Kaggle aggregate has a near duplicate in another public
corpus. Removing those duplicates lowers cross-corpus F1 by up to 30 points for
a single pair, while removing the same number of emails at random changes
nothing. On decontaminated, unseen data a small open LLM used zero-shot is as
good as a trained classical model and much better on AI-written phishing, for a
few cents per thousand emails, though a current commercial model is still
ahead of all of them. Which model is best depends on the operating point: at a
realistic phishing rate the false alarm rate decides, and a free classical
filter in front of a small LLM beats either part alone on both cost and false
alarms. Finally, the corpora themselves deserve more scepticism than they
usually get: four fifths of their positive emails are spam rather than
phishing, and a corpus feature that looks like a shortcut needs a paired test
before it is believed, in either direction.
High accuracy on one dataset is easy. Honest evaluation is the hard part. All
code, raw model answers and results are public so that others can check and
extend this work.

\section*{Acknowledgment}
We thank our Computer Security course instructor at United International
University for the feedback on our literature review and proposal.
